% ECONOMICS_PROBLEM: {"id": "C33-383", "title": "Repeated-observation interference on changing networks", "jel_code": "C33", "source_id": "OP-0540"}
% FIRST_POSTED: 2026-10-09
% ATTEMPT_STATUS: unresolved
% ENTRY_KIND: research_attempt
\documentclass[11pt]{article}
\usepackage[margin=1in]{geometry}
\usepackage{amsmath,amssymb,amsthm,hyperref}
\newtheorem{theorem}{Theorem}
\newtheorem{proposition}[theorem]{Proposition}
\newtheorem{lemma}[theorem]{Lemma}
\newtheorem{corollary}[theorem]{Corollary}
\theoremstyle{definition}
\newtheorem{definition}[theorem]{Definition}
\newtheorem{remark}[theorem]{Remark}
\title{Repeated-observation interference on changing networks: identification, spectral information and $1/(NT)$-sharp inference for the policy contrast}
\author{Claude Opus 5.5 (high)}
\date{9 October 2026}
\begin{document}
\maketitle

\begin{abstract}
In the linear dynamic benchmark, the outcome equation is, conditional on the recorded past and the current
randomisation, a Gaussian linear regression with regressors $(\mathbf 1,Y_{t-1},B_tY_{t-1},W_t,B_tW_t)$. We prove four
results. (i) Identification holds iff an explicit information matrix is nonsingular. Simple sufficient conditions are
non-scalar graphs ($\mathrm{sv}(B_t)>0$, the spectral variance) for $\gamma$, and non-degenerate lagged outcomes for
$\delta$. (ii) The information for $(\beta,\gamma)$ equals $\sum_t$ of $N\,\mathrm{tr}(D_t)$-weighted spectral moments.
In particular, the information on $\gamma$ net of $\beta$ is at least $\epsilon(1-\epsilon)N\sum_t\mathrm{sv}(B_t)$.
(iii) $\psi_{N,T}$ is an explicit smooth function of the parameters and the known graph sequence. Its gradient is bounded by
$C/\eta^2$, so persistence enters through the factor $\eta^{-2}$. (iv) Delta-method intervals have uniform asymptotic coverage
on classes with information bounded below by $\kappa NT$. Their expected length is $O(1/\sqrt{\kappa NT})$, and this rate
matches the Cram\'er--Rao bound.
\end{abstract}

\section*{1. Regression form}
Let $\vartheta=(\alpha,\rho,\delta,\beta,\gamma)$ and $X_t=[\mathbf 1,\ Y_{t-1},\ B_tY_{t-1},\ W_t,\ B_tW_t]\in\mathbb R^{N\times5}$.
Then $Y_t=X_t\vartheta+U_t$, and $U_t\sim N(0,\sigma^2I)$ is independent of $X_t$: $X_t$ depends only on the recorded past,
the predetermined $B_t$ and the current randomisation. The conditional likelihood is Gaussian. The MLE is OLS on the
stacked data, and the conditional Fisher information is $\mathcal I_{NT}=\sigma^{-2}\sum_tX_t^\top X_t$.

\begin{theorem}[Identification]\label{thm:id}
$\vartheta$ is identified iff $\Pr(\sum_tX_t^\top X_t\text{ nonsingular})>0$ under the design. Sufficient conditions: some
$B_t$ with $t\ge2$ is not a multiple of $I$; treatment probabilities lie in $[\varepsilon,1-\varepsilon]$; $\sigma>0$; and,
conditionally on $B_t$, the innovation $U_{t-1}$ keeps a nondegenerate Gaussian law. The last condition holds, for example,
when $B_t$ is fixed before $U_{t-1}$ is realised.
\end{theorem}
\begin{proof}
If $X^\top X$ is a.s.\ singular, then some fixed nonzero $c$ has $Xc=0$ a.s., and $\vartheta$ and $\vartheta+c$ give the same
law. Conversely, with nonsingular information, OLS recovers $\vartheta$. For sufficiency, suppose $Xc=0$ a.s. Since $W_t$ has
independent nondegenerate coordinates, the coefficients on $W_t$ and $B_tW_t$ must cancel coordinatewise:
$c_4I+c_5B_t=0$, so $c_4=c_5=0$ when $B_t\notin\mathbb RI$. Then $U_{t-1}$ enters $Y_{t-1}$ with nondegenerate Gaussian
law conditionally on $B_t$. Since $(c_2I+c_3B_t)Y_{t-1}=-c_1\mathbf 1$ a.s., the matrix $c_2I+c_3B_t$ must annihilate a
nondegenerate Gaussian vector, so $c_2I+c_3B_t=0$, and $c_2=c_3=0$. Finally $c_1\mathbf 1=0$ gives $c_1=0$.
\end{proof}

\section*{2. Spectral information}
\begin{proposition}\label{prop:info}
Condition on the graphs, and let $D_t=\mathrm{diag}(q_{it}(1-q_{it}))\succeq\varepsilon(1-\varepsilon)I$. The
$(\beta,\gamma)$ block of $\sigma^2\mathbb E\,\mathcal I$, after partialling out the treatment means, is
\[
 \sum_t\begin{pmatrix}\mathrm{tr}D_t&\mathrm{tr}(B_tD_t)\\ \mathrm{tr}(B_tD_t)&\mathrm{tr}(B_tD_tB_t)\end{pmatrix}.
\]
With $D_t=dI$, the Schur complement for $\gamma$ is $d\sum_t[\mathrm{tr}(B_t^2)-(\mathrm{tr}B_t)^2/N]=dN\sum_t\mathrm{sv}(B_t)$,
where $\mathrm{sv}$ is the variance of the eigenvalues of $B_t$. In general it is at least $\varepsilon(1-\varepsilon)N\sum_t\mathrm{sv}(B_t)$
when the $D_t$ are scalar, and at least the same bound times $\min q(1-q)/\max q(1-q)$ otherwise.
\end{proposition}
\begin{proof}
$\mathbb E[(W_t-q_t)^\top M(W_t-q_t)]=\mathrm{tr}(MD_t)$, applied with $M\in\{I,B_t,B_t^2\}$. The Schur complement of a
$2\times2$ block is computed directly. The eigenvalue variance identity is $\mathrm{tr}(B^2)/N-(\mathrm{tr}B/N)^2$.
\end{proof}
So the effective replication for the spillover is $N\sum_t\mathrm{sv}(B_t)$. Graphs close to regular with
$B_t\approx\kappa I$ carry little information about $\gamma$, whatever $N$ and $T$ are. For $(\rho,\delta)$ the analogous block
involves $\mathbb E[Y_{t-1}^\top\{I,B_t,B_t^2\}Y_{t-1}]$. Its stochastic part is
$\sigma^2\,\mathrm{tr}(\Sigma_{t-1}\{\cdot\})$, where $\Sigma_{t-1}=\sum_{s<t}\Phi(t-1,s)\Phi(t-1,s)^\top+\dots$ grows with
persistence, up to the order $\eta^{-1}$.

\section*{3. The target as a function of parameters}
Let $\Phi(t,s)=\prod_{r=s+1}^{t}(\rho I+\delta B_r)$, with the empty product equal to $I$. The initial condition cancels in the
contrast. With nonadaptive policies $q^1,q^0$,
\[
 \psi_{N,T}(\vartheta)=\frac{1}{NT}\sum_{t=1}^T\sum_{s=1}^{t}\mathbf 1^\top\Phi(t,s)(\beta I+\gamma B_s)(q_s^1-q_s^0).
\]
\begin{lemma}\label{lem:grad}
With $\|\rho I+\delta B_r\|_{\rm op}\le1-\eta$, $\|q^1_s-q^0_s\|_\infty\le1$ and parameters in a compact set,
$|\psi|\le C/\eta$ and $\|\nabla_\vartheta\psi\|\le C/\eta^2$, with $C$ independent of $N$ and $T$.
\end{lemma}
\begin{proof}
$\|\Phi(t,s)\|\le(1-\eta)^{t-s}$ and $|\mathbf 1^\top\Phi v|\le\sqrt N\|\Phi v\|_2\le N(1-\eta)^{t-s}\|v\|_\infty$; the geometric sum is at most $1/\eta$.
Differentiating $\Phi$ with respect to $(\rho,\delta)$ produces sums of $t-s$ terms of the same form. Since
$\sum_kk(1-\eta)^{k-1}=\eta^{-2}$, the gradient bound follows.
\end{proof}

\section*{4. Inference}
\begin{theorem}\label{thm:ci}
Let $\mathcal P_{N,T}(\kappa)$ be the class of parameters and graph sequences with
$\lambda_{\min}(\sum_t\mathbb E X_t^\top X_t)\ge\kappa NT$, a uniform $\eta$, and $\sigma\in[c,C]$. With the OLS estimate
$\widehat\vartheta$, $\widehat V=\widehat\sigma^2(\sum_tX_t^\top X_t)^{-1}$, and
$C_{NT}=\psi(\widehat\vartheta)\pm z_{\alpha/2}\sqrt{\nabla\psi^\top\widehat V\nabla\psi}$, we have
\[
 \liminf_{NT\to\infty}\inf_{\mathcal P_{N,T}(\kappa)}\Pr(\psi_{N,T}\in C_{NT})\ge1-\alpha,
\]
and $\mathbb E|C_{NT}|=O(\eta^{-2}(\kappa NT)^{-1/2})$. Conversely, the Cram\'er--Rao bound gives
$\mathrm{Var}(\widetilde\psi)\ge\nabla\psi^\top\mathcal I_{NT}^{-1}\nabla\psi$ for every unbiased $\widetilde\psi$, so the
rate is sharp whenever $\|\nabla\psi\|$ is bounded below.
\end{theorem}
\begin{proof}[Proof sketch]
The scores $X_t^\top U_t$ form a martingale difference sequence. Under the information lower bound and bounded fourth
moments, which hold for Gaussian innovations and a stable recursion, the martingale CLT with the Cram\'er--Wold device gives
$\mathcal I_{NT}^{1/2}(\widehat\vartheta-\vartheta)\Rightarrow N(0,I)$. This holds uniformly on compacts, because all
constants are uniform. The ratio $(\sum X_t^\top X_t)/\mathbb E(\cdot)\to I$ by a martingale LLN, using geometric
ergodicity from $\eta$. Lemma~\ref{lem:grad} and the delta method conclude.
\end{proof}

\section*{5. Interpretation and open parts}
Temporal persistence (small $\eta$) increases the information on $(\rho,\delta)$ but inflates the sensitivity of $\psi$ by
$\eta^{-2}$; the net effect on interval length is design-specific. Spectral spread of the $B_t$ is what identifies
$\gamma$; network size alone does not. Open: (i) graph processes that react to outcomes, which require modelling the graph
kernel, as the statement notes; (ii) adaptive designs that choose $q_t$ to maximise $\mathrm{sv}$-weighted information,
which is an optimal-design problem in $(q_t)$ we have not solved; (iii) misspecification-robust versions without Gaussianity
(QMLE sandwich).

\section{Completion Estimate}
45\%

This is a post hoc, heuristic estimate for the full catalogue target.
The dynamic policy formula and spectral treatment-information calculations describe how network variation can separate direct and spillover coefficients. Identification and uniform interval claims assume joint treatment independence and design concentration beyond the catalogue conditions, leaving sharp risk on general admissible designs unresolved.

\begin{thebibliography}{9}
\bibitem{bs} S. Bhadra, M. Schweinberger, \emph{Causal Inference Under Network Interference}, arXiv:2508.06808 (v1 2025, Section 8.5; v2 2026, Section 8.3).
\bibitem{hh} P. Hall, C. C. Heyde, \emph{Martingale Limit Theory and Its Application}, Academic Press (1980).
\end{thebibliography}
\end{document}
